\input{../tex-inputs/latex-leseansicht-vorspann}

\section[Arthur Schnitzler und Hugo von Hofmannsthal an Gustav Schwarzkopf, 15. 3. 1905]{L04392 Arthur Schnitzler und Hugo von Hofmannsthal an Gustav Schwarzkopf, 15. 3. 1905}
\nopagebreak\mylabel{L04392v}
\rehead{ }\normalsize\beginnumbering\briefempfaengerindex{Schwarzkopf, Gustav@\textsc{Schwarzkopf, Gustav}!zzzHofmannsthal, Hugo von@\emph{von Hugo von Hofmannsthal}!1905-03-152@{15. 3. 1905}|(be}\briefempfaengerindex{Schwarzkopf, Gustav@\textsc{Schwarzkopf, Gustav}!zzzSchnitzler, Arthur@\emph{von Arthur Schnitzler}!1905-03-152@{15. 3. 1905}|(be}
\toendnotes[C]{\smallbreak\pagebreak[2]}
\correspDesc{Versand  durch Arthur Schnitzler, Hugo von Hofmannsthal am 15. 3. 1905 in Dubrovnik
\newline{}Übermittlung  am 15. 3. 1905 in Dubrovnik
\newline{}Zustellung  am 17. 3. 1905 in Wien
\newline{}Erhalt  durch Gustav Schwarzkopf am 17. 3. 1905 in Wien}\toendnotes[C]{\smallbreak}
\Standort{DLA, A:Schnitzler, HS.1985.1.1897.}
\physDesc{Bildpostkarte, 101 Zeichen
\newline{}Handschrift Arthur Schnitzler: Bleistift, lateinische Kurrent
\newline{}Handschrift Hugo von Hofmannsthal: Bleistift, lateinische Kurrent
\newline{}Versand: 1) Stempel: »\nobreak{}\oindex{Gruž@\textbf{Gruž}|pwk}Gruž Gra{[}vosa{]}, 15/3 0{[}5{]}\nobreak{}«.   2) Stempel: »\nobreak{}\oindex{I., Innere Stadt@\textbf{I., Innere Stadt}|pwk}Wien 1/1 1, 17/3. 05, IX, Bestellt\nobreak{}«. }\pstart{}{\pb}Herrn Gustav Schwarzkopf\pend{}\pstart{}Wien I\oindex{I., Innere Stadt@\textbf{I., Innere Stadt}|pw}\pend{}\pstart{}Tiefer Graben 23\oindex{Wien@\textbf{Wien}!I., Innere Stadt@\textbf{I., Innere Stadt}!Tiefer Graben 23@\textbf{Tiefer Graben 23}|pw}\pend{}{\bigskip}
\pstart
           \centering{}{\pb}\textcolor{gray}{\textbf{Hamburg–Amerika Linie\orgindex{Hamburg-Amerika-Linie@Hamburg-Amerika-Linie|pw}.}}\pend
           
\pstart
           \centering{}\textcolor{gray}{\textbf{Am Bord des Doppelschrauben-Dampfers »METEOR\orgindex{Meteor@Meteor|pw}«.}}\pend
           \vspace{1em}
\pstart
           {\pb}Ragusa\oindex{Dubrovnik@\textbf{Dubrovnik}|pw}{ }5/3 905\pend
           \vspace{0.5em}
\pstart
           Herzliche Grüße!\pend
           \pstart Ihr \spacefill\mbox{A. S.}\pend{}\selectlanguage{ngerman}\vspace{1em}
\pstart
           \noindent{}{[}hs. Hofmannsthal:{]} Viele Grüße!\pend
           \pstart \spacefill\mbox{Hugo.}\pend{}\selectlanguage{ngerman}\endnumbering\briefempfaengerindex{Schwarzkopf, Gustav@\textsc{Schwarzkopf, Gustav}!zzzHofmannsthal, Hugo von@\emph{von Hugo von Hofmannsthal}!1905-03-152@{15. 3. 1905}|)be}\briefempfaengerindex{Schwarzkopf, Gustav@\textsc{Schwarzkopf, Gustav}!zzzSchnitzler, Arthur@\emph{von Arthur Schnitzler}!1905-03-152@{15. 3. 1905}|)be}\mylabel{L04392h}
\begin{anhang}
\end{anhang}\newcommand{\dateiname}{L04392}\newcommand{\titel}{Arthur Schnitzler und Hugo von Hofmannsthal an Gustav Schwarzkopf, 15. 3. 1905}\newcommand{\editorInnen}{Herausgegeben von Jahnke, SelmaMüller, Martin Anton}\input{../tex-inputs/latex-leseansicht-abspann}
